\section{Additional Method and Protocol Details}
\label{app:method}

VERA assesses semantic browser actions submitted to the executor, not every
atomic mouse or keyboard event used internally to realize them. Ordinary
verification links one such action to a frozen claim and outcome evidence.
Replay consists of previously observed semantic actions; every replay step
receives its own risk record and is counted separately from ordinary candidate
attempts. If risk inference fails, execution continues and the error is retained
in the audit record. This behavior keeps risk recognition observational in the
evaluated system.

Formal RQ1/RQ2 runs start from designated entry URLs in clean browser sessions
with resettable state restored. Each application--condition pair has three runs
and a maximum of 25 ordinary candidate attempts. Configurations record model
settings, viewport, seeds, retry rules, stopping conditions, and replay limits.

\section{Prompts, Schemas, and Risk Taxonomy}

The artifact package includes the functional-claim and evidence schemas,
versioned proposal and judgment prompts, and
\texttt{risk\_taxonomy\_v1.json}. Risk outputs contain a binary
\texttt{potential\_risk}, one primary \texttt{risk\_type} when applicable, and
brief GUI-grounded evidence. The same taxonomy is supplied across compared C3
conditions and therefore serves as shared output vocabulary.

\section{Annotation and Uncertainty Protocols}

Initial C1--C3 annotations were generated with AI assistance and then reviewed
and revised item by item by a human author under frozen guides. The resulting
labels are used as gold; the process is not independent double annotation. C1
gold is aggregated by application, semantic location, and canonical action
identifier. C2 uses frozen per-application coverage inventories and a reviewed
claim-to-inventory ledger. C3 permits multiple acceptable risk types for
compound external cases while the assessor returns one primary type.

External C3 uncertainty estimates use 10,000 cluster-bootstrap samples, with
each candidate context pair treated as one cluster and each diversity case as a
singleton (seed 20260920). Differences whose intervals cross zero are reported
as unresolved point-estimate directions.

\section{Result Artifacts and Additional Cases}

Machine-readable metrics, reviewed gold, artifact manifests, and complete
reports are organized under \texttt{paper/experiments/results/E002--E004}.
These packages retain failed executions, incomplete evidence, and negative
examples rather than filtering them from analysis.

The C1 false admission concerns a form-completion claim for which grouped
evidence did not directly show that all required fields were completed. The
missed supported claim concerns product sorting despite a visible ordering
change. In external C3, a recurring binary error treats entry into a
consequential workflow as if the immediate action completed the downstream
consequence; context nevertheless corrects substantially more risk-type
decisions than it introduces.
